\section{Einleitung}

Magnetische Effekte begegnen in der Physik häufig zunächst in vereinfachter Form, etwa als ideales Feld einer Spule oder als lineare Beziehung zwischen Feldstärke und Magnetisierung. In realen Materialien zeigt sich jedoch ein deutlich komplexeres Verhalten. Insbesondere ferromagnetische Stoffe reagieren nicht nur auf ein aktuelles Magnetfeld, sondern behalten eine „Erinnerung“ an dessen Verlauf. Dies ist bekannt als Hysterese und ist sowohl aus physikalischer Sicht von Bedeutung als auch für zahlreiche technische Anwendungen entscheidend. \\
Ziel des vorliegenden Versuchs ist es, diese Zusammenhänge experimentell zu untersuchen. Ausgangspunkt bildet die Messung des Magnetfeldes einer stromdurchflossenen Spule, wodurch die räumliche Verteilung magnetischer Felder dargestellt wird.  \\
Darauf aufbauend wird das Verhalten ferromagnetischer Materialien im äußeren Magnetfeld betrachtet. Anhand von Hysteresekurven lässt sich nachvollziehen, wie stark die Magnetisierung vom bisherigen Feldverlauf abhängt und welche charakteristischen Kenngrößen das Materialverhalten bestimmen. Ein besonderer Fokus liegt dabei auf dem Vergleich eines Vollkerns mit einem geblätterten Eisenkern, da sich hier grundlegende Unterschiede hinsichtlich Energieverlusten und magnetischer Eigenschaften zeigen. \\
Ergänzend dazu wird die induktive Kopplung zweier Spulen genutzt, um materialabhängige Größen wie die relative Permeabilität zu bestimmen. Diese beschreibt, wie effektiv ein Material magnetische Felder verstärken und leiten kann und ist damit eine zentrale Kenngröße für technische Anwendungen. \\
Der Versuch verbindet somit grundlegende physikalische Konzepte mit praxisnahen Fragestellungen. Insbesondere wird deutlich, welche Rolle die Materialwahl für die Effizienz elektromagnetischer Systeme spielt und warum bestimmte Bauformen, wie geblätterte Kerne, in der Technik bevorzugt eingesetzt werden.

